\subsubsection{Random pile sizes: the lower bound}

\begin{corollary}[Random cut arrays: lower bound]
\label{cor:lower-random}
Fix the parameters of Theorem~\ref{thm:empirical-lower}, and let
$\mathcal B_N$ be the set of deterministic cut arrays satisfying
either of its two alternatives. For any probability law $Q_N$ on
cut arrays of length $K=\lfloor T\log N\rfloor$, put
\[
 \overline P_N^K=\int P_N^K(\mathbf n)\,Q_N(d\mathbf n).
\]
Then, uniformly in $Q_N$,
\[
 \dTV(\overline P_N^K,\nuUnif)\ge Q_N(\mathcal B_N)-o(1).
\]
In particular, the annealed distance tends to one if
$Q_N(\mathcal B_N)\to1$. The cut arrays may have arbitrary
dependence between their columns; conditional on the entire array,
the shuffle law in this statement uses independent uniform
interleavings. The same conclusion holds for the corresponding
chronological forward shuffle by inversion.
\end{corollary}

\begin{proof}
Set
\[
 e=\min\{\eta_1/100000,1/1000\},\qquad
 b=\min\{\eta_0,e\},\qquad
 m_N=\lfloor N^{b/2}\rfloor.
\]
Let $\mathcal R_N$ be the event that the permutation has an
increasing run of length $m_N$. For sufficiently large $N$,
$2\le m_N\le N^b/2$. Proposition~\ref{prop:largest-prefix} and
its proof therefore give, uniformly over arrays with
$\beta^{(N)}(S_N)\ge b$,
\[
 P_N^K(\mathbf n)(\mathcal R_N)
 \ge1-\exp\left(-\frac{1-\log2}{2}N^b\right),\qquad
 \nuUnif(\mathcal R_N)\le\frac{N}{m_N!}=o(1).
\]

Let $\mathcal C_N$ consist of the arrays satisfying the second
alternative of Theorem~\ref{thm:empirical-lower} and
$\beta^{(N)}(S_N)<e$. If the cold parameters are nonpositive or
violate the restriction in that theorem, this set is empty.
For each $\mathbf n\in\mathcal C_N$, fix one cold set
$H_{\mathrm{rel},\mathbf n}$ from
Lemma~\ref{lem:cold-geometry}, and write
\[
 M_{\mathbf n}=|H_{\mathrm{rel},\mathbf n}|,\qquad
 H_{q,\mathbf n}=(q+H_{\mathrm{rel},\mathbf n})\cap[N],\qquad
 \delta=\eta_1/5000.
\]
Define the permutation event
\[
 \mathcal A_{\mathbf n}
 =\left\{\sigma:\exists q\in\{0,\ldots,N\},\quad
 X_{H_{q,\mathbf n}}(\sigma)>
 \frac{|E(H_{q,\mathbf n})|}{2}
       +M_{\mathbf n}^{1/2+\delta/2}\right\}.
\]
Lemma~\ref{lem:many-edges}, followed by the conditional comparison
in the proof of Lemma~\ref{lem:lalley}, gives
\[
 P_N^K(\mathbf n)(\mathcal A_{\mathbf n})=1-o(1)
\]
uniformly over $\mathcal C_N$. Indeed, on the event supplied by
Lemma~\ref{lem:many-edges}, the conditional failure probability
is at most $\exp(-M_{\mathbf n}^{2\delta}/144)$.
On the other hand, $M_{\mathbf n}\ge N^{\eta_1/3000}$, and
Lemma~\ref{lem:uniform-ascent} gives
\[
 \nuUnif(\mathcal A_{\mathbf n})
 \le (N+1)\exp(-M_{\mathbf n}^{\delta})
 \le (N+1)\exp\left(-N^{\eta_1^2/15000000}\right).
\]

There are at most $(N+1)^{kK}$ possible cut arrays. Thus the
single event
\[
 \mathcal E_N=\mathcal R_N\cup
           \bigcup_{\mathbf n\in\mathcal C_N}\mathcal A_{\mathbf n}
\]
has uniform-law probability bounded by
\[
 \nuUnif(\mathcal E_N)
 \le\frac{N}{m_N!}
 +(N+1)^{kK+1}\exp\left(-N^{\eta_1^2/15000000}\right)
 =o(1),
\]
since $K\le T\log N$.
Every array in $\mathcal B_N$ either has
$\beta^{(N)}(S_N)\ge b$ or belongs to $\mathcal C_N$.
Consequently
$P_N^K(\mathbf n)(\mathcal E_N)\ge1-o(1)$ uniformly on
$\mathcal B_N$. Integrating this common-event bound gives
\[
 \overline P_N^K(\mathcal E_N)\ge Q_N(\mathcal B_N)-o(1).
\]
Subtracting $\nuUnif(\mathcal E_N)$ proves the asserted
total-variation lower bound. All errors depend only on the fixed
parameters, so the conclusion is uniform in $Q_N$.
\end{proof}
